\section{Elección de parámetros y ventanas comunes}
\label{sec:parametros}

Las estimaciones de las notas anteriores se convierten en una cota para
$H_\psi$ después de dividir el exceso aritmético por el coeficiente que
produce la selección de caracteres. Esa división cambia el tamaño de
cada error. Además, el criterio analítico utiliza un intervalo entero:
los datos que permiten acotar el exceso deben existir de manera compatible
en todos sus puntos. Esta nota desarrolla el cálculo de los exponentes
y la coordinación de las elecciones.

\fuente{A, §§10--11, pp.~46--52, con la extracción de la Proposición~4.3,
pp.~16--17; B, §6, Lema~6.1 y Proposición~6.3, pp.~64--67.
Se mantienen los ejemplares A de 55 páginas y B de 73 páginas del
21 de noviembre. Las deducciones numéricas y lógicas se prueban aquí;
la realización aritmética de todas sus hipótesis se identifica como
entrada de la propuesta.}

\subsection{La estimación que debe producir la construcción}

Sea $\psi$ un carácter primitivo de orden $o$. Conservamos la notación
de la Nota~\ref{sec:caracteres}:
\begin{equation}
 H_\psi(x)=G_1(x)+\frac{1}{3o+1}
       \sum_{j=1}^{o-1}\Ree G_{\psi^{[j]}}(x).
 \label{eq:par-hpsi}
\end{equation}
Aquí $\psi^{[j]}$ es el carácter primitivo que induce la potencia
$\psi^j$. Para $o=1$ la suma es vacía y $H_\psi=G_1$.
B llama $\mathcal G_\psi$ a esta función. Los coeficientes de
\eqref{eq:par-hpsi} son positivos; los valores de la función pueden
tener ambos signos.

La meta es la propiedad
\begin{equation}
 \forall\eta>0\ \forall A\ge3\ \exists Y\ge A\
 \forall x\in[Y,Y^{3/2}],\qquad H_\psi(x)\le x^{1/2+\eta}.
 \label{eq:par-meta}
\end{equation}
El carácter $\psi$ se fija antes de estas elecciones. Los cuerpos,
pesos y mapas auxiliares podrán variar con $x$ si las constantes de la
estimación y el intervalo ya están fijados. La función $H_\psi$ permanece
la misma en todos los puntos.

La Nota~\ref{sec:ventanas}, Proposición~\ref{vent-corolario}, explica por
qué \eqref{eq:par-meta}, junto con las hipótesis analíticas allí
comprobadas, excluye los ceros de $L(s,\psi)$ con parte real mayor
que $1/2$. El trabajo aritmético consiste en producir esta cota
unilateral en ventanas arbitrariamente lejanas.

\subsection{El presupuesto se forma después de cancelar}

La comparación de pendientes de la Proposición~\ref{prop:periodos-balance},
bajo sus hipótesis, entrega un exceso de la forma
\[
 \Delta=\sum_{j,t}
   \bigl(\mathcal A(a_t,t,w_j)-\mathcal M(a_t,w_j)\bigr).
\]
La cantidad $\mathcal A$ reúne ganancias de divisibilidad; $\mathcal M$
es el término principal que contiene los logaritmos y la derivada
logarítmica que interviene en la normalización del periodo
arquimediano. En A, Proposición~10.1, la comparación se organiza,
para una traslación fijada, mediante
\begin{equation}
 \Delta\le M_0+E_{\mathrm{HN}}+
       \sum_i p_iL_i+E_D+E_C+A_\Sigma E_T+E_P,
 \qquad A_\Sigma=\sum_t a_t,\quad \sum_i p_i=1.
 \label{eq:par-costes}
\end{equation}
Los pesos $p_i$ son positivos. $M_0$ controla la pendiente inicial;
$E_{\mathrm{HN}}$ es la pérdida del marco; $L_i$ recoge los costes
integrales; los demás términos registran las comparaciones de altura,
nivel, cardinalidad y ganancia finita. El subíndice de $A_\Sigma$ evita
confundir esta suma con la fuente A o con el umbral de
\eqref{eq:par-meta}.

La formación de \eqref{eq:par-costes} usa cancelaciones exactas ya
preparadas en la Nota~\ref{sec:periodos}. El retorcimiento finito común
aparece con signos opuestos en la fuente y el destino. El término de
nivel $mA_\Sigma\log g_0$ se cancela con el que aporta el conteo de
palabras. Finalmente, la identidad normalizada
\begin{equation}
 \log\Omega=\gamma_E+\log(\pi/v)-\Ree\ell_\kappa
 \label{eq:par-periodo}
\end{equation}
cancela los logaritmos de $v/\pi$ procedentes de los pesos
arquimedianos. Aquí $\kappa$ es el carácter cuadrático del periodo y
$\ell_\kappa=L'(1,\kappa)/L(1,\kappa)$.
La cota \eqref{eq:par-costes} se aplica al resto que queda tras estas
cancelaciones. Sustituir prematuramente los términos firmados por sus
valores absolutos perdería las cancelaciones que identifican
$\mathcal M$.

Las Notas~\ref{sec:geometria} y~\ref{sec:cohomologia} explican cómo
varias entradas geométricas y cohomológicas controlan esos restos.
En este bloque utilizamos como hipótesis que su suma queda acotada por
$x^{C\varepsilon}\mathcal B$, con un exponente $C$ fijo.
La expresión $\mathcal B$ se analiza a continuación. La estimación
debe valer para cada traslación que usa la extracción de caracteres:
\[
 \Delta_u(x)\le E(x)\qquad\text{para todo }u\in A_0.
\]
El sistema de datos, el presupuesto y los factores normalizadores
son comunes, con las dependencias uniformes de esa extracción.

\subsection{Parámetros, densidades y coeficiente de extracción}

Fijemos $x\ge3$, $\varepsilon>0$ y $C\ge0$. Escribimos
$n$ para el tamaño del factor abeliano del grupo auxiliar,
$t=|T_1|$ para el número de posiciones superiores y
$v_2=|T_2|$ para el tamaño del conjunto grueso. En A estos dos últimos
números se escriben $r_T$ y $n_2$; $n^2$ denotará siempre el
cuadrado de $n$.

La longitud de la torre theta es $m$. El número $k$ es la dimensión
inicial y $b>0$ la escala de los pesos grandes. Usamos $a_E\ge0$
para una cota superior de las amplitudes $a_t$. El cálculo requiere
\begin{equation}
 \begin{gathered}
 0\le a_E\le2xb,\qquad
 t\ge\frac n2x^{-\varepsilon/8},\qquad
 nx^{-\varepsilon}\le v_2\le n,\\
 x^{1/2-\varepsilon}\le m\le x^{1/2+\varepsilon/2},
 \qquad
 xm\le k\le n^2x^{1/2-\varepsilon}.
 \end{gathered}
 \label{eq:par-rangos}
\end{equation}
Los números $n,t,v_2,m$ son enteros positivos en la aplicación.
La estimación elemental sigue siendo válida para valores reales
positivos de esos parámetros y de $k$. Relajar el dominio en un
lema numérico no construye los espacios y conjuntos que los realizan.

Definamos
\begin{equation}
 Q=tmb>0,\qquad d=\frac{v_2}{n},\qquad
 \mathcal B=
 ta_E\left(1+\frac{m^2}{k}\right)
 +m\left(2nmbx^\varepsilon+
               \frac{v_2kbx^{2\varepsilon}}{n^2}\right).
 \label{eq:par-presupuesto}
\end{equation}
La letra $d$ designa aquí una densidad; en el apartado de ventanas
comunes volveremos a escribir explícitamente los primos auxiliares
como $q_{\mathrm{aux}},d_{\mathrm{aux}}$.

El coeficiente que acompaña a la combinación retenida es $Qd/2$.
El factor $d/2$ procede de \eqref{car-factor-medio}; $t$ cuenta las
posiciones y $m$ las etapas. Por ello, una cota grande para el exceso
puede dar una cota pequeña para $H_\psi$ si se compara con este
coeficiente, conservando las pérdidas debidas a las densidades.

\begin{center}
\small
\begin{tabularx}{\textwidth}{@{}p{0.18\textwidth}YY@{}}
\toprule
\textbf{Cantidad} & \textbf{Cota utilizada} & \textbf{Efecto al normalizar}\\
\midrule
$a_E/b$ & $a_E/b\le2x$ & Produce los términos $x/m$ y $xm/k$.\\
$t/n$ & $t/n\ge x^{-\varepsilon/8}/2$ &
El factor $n/t$ cuesta a lo sumo $2x^{\varepsilon/8}$.\\
$v_2/n=d$ & $d\ge x^{-\varepsilon}$ &
Despejar $H_\psi$ cuesta a lo sumo $x^\varepsilon$.\\
$m$ & Entre $x^{1/2-\varepsilon}$ y $x^{1/2+\varepsilon/2}$ &
Mantiene próximos a $\sqrt{x}$ los costes opuestos.\\
$k$ & $xm\le k\le n^2x^{1/2-\varepsilon}$ &
Controla simultáneamente un término inverso y otro directo en $k$.\\
\bottomrule
\end{tabularx}
\end{center}

\subsection{Por qué la longitud de torre está cerca de la raíz}

El primer término de $\mathcal B/Q$ es $a_E/(mb)\le2x/m$.
El tercero es $2(n/t)mx^\varepsilon$. Salvo los factores
$x^{O(\varepsilon)}$ permitidos por \eqref{eq:par-rangos},
uno disminuye y el otro aumenta al crecer $m$.

\begin{lema}[Equilibrio de dos costes]
\label{lem:par-equilibrio}
Para $x,m>0$ se tiene
\[
 \frac xm+m\ge2\sqrt{x},
\]
con igualdad si y solo si $m=\sqrt{x}$. Si $m$ debe ser entero,
la elección $m=\lceil\sqrt{x}\rceil$ satisface el rango de $m$ en
\eqref{eq:par-rangos} siempre que $x\ge2^{2/\varepsilon}$.
\end{lema}

\begin{proof}
La identidad
\[
 \frac xm+m-2\sqrt{x}
   =\left(\sqrt m-\sqrt{x/m}\right)^2
\]
prueba la desigualdad y su caso de igualdad. Para $x\ge1$,
\[
 \sqrt{x}\le\lceil\sqrt{x}\rceil
       \le\sqrt{x}+1\le2\sqrt{x}.
\]
La cota inferior de \eqref{eq:par-rangos} se sigue de
$x^{-\varepsilon}\le1$; la superior se sigue de
$2\le x^{\varepsilon/2}$.
\end{proof}

El rango de $k$ impone otra compatibilidad:
\begin{equation}
 xm\le n^2x^{1/2-\varepsilon}
 \quad\Longleftrightarrow\quad
 n\ge m^{1/2}x^{1/4+\varepsilon/2}.
 \label{eq:par-compatibilidad}
\end{equation}
Si se usa también $m\ge x^{1/2-\varepsilon}$, resulta la condición
necesaria $n\ge\sqrt{x}$. Elegir $m$ por sí solo deja abierta la
construcción del grupo auxiliar, de la dimensión y de las posiciones
que cumplen todos los rangos. Esa realización simultánea será una
hipótesis explícita.

\subsection{La cuenta de los cuatro términos}

\begin{lema}[Presupuesto normalizado]
\label{lem:par-presupuesto}
Bajo \eqref{eq:par-rangos}, la cantidad definida en
\eqref{eq:par-presupuesto} satisface
\begin{equation}
 \frac{\mathcal B}{Q}\le10x^{1/2+13\varepsilon/8}.
 \label{eq:par-diez}
\end{equation}
En consecuencia, si $E\le x^{C\varepsilon}\mathcal B$, entonces
\begin{equation}
 \frac{2E}{Qd}\le20x^{1/2+(C+21/8)\varepsilon}.
 \label{eq:par-veinte}
\end{equation}
\end{lema}

\begin{proof}
La división por $Q=tmb$ da la igualdad
\begin{equation}
 \frac{\mathcal B}{Q}
 =\frac{a_E}{mb}+\frac{a_Em}{bk}
     +\frac{2nmx^\varepsilon}{t}
     +\frac{v_2kx^{2\varepsilon}}{n^2t}.
 \label{eq:par-cuatro}
\end{equation}
Las cuatro cotas se obtienen por separado. Primero,
\[
 \frac{a_E}{mb}\le\frac{2x}{m}\le2x^{1/2+\varepsilon},
 \qquad
 \frac{a_Em}{bk}\le\frac{2xm}{k}\le2.
\]
La densidad de $T_1$ y la cota superior de $m$ dan
\[
 \frac{2nmx^\varepsilon}{t}
 \le4m x^{9\varepsilon/8}
 \le4x^{1/2+13\varepsilon/8}.
\]
Para el último término usamos primero $v_2\le n$, después
$n/t\le2x^{\varepsilon/8}$ y finalmente la cota superior de $k$:
\[
 \frac{v_2kx^{2\varepsilon}}{n^2t}
 \le\frac{kx^{2\varepsilon}}{nt}
 \le\frac{2k}{n^2}x^{17\varepsilon/8}
 \le2x^{1/2+9\varepsilon/8}.
\]
Cada potencia obtenida es a lo sumo
$x^{1/2+13\varepsilon/8}$, pues $x\ge1$ y $\varepsilon>0$.
Sumar los coeficientes $2+2+4+2=10$ prueba
\eqref{eq:par-diez}. La desigualdad
$d^{-1}\le x^\varepsilon$ convierte esa cota en
\eqref{eq:par-veinte}.
\end{proof}

El incremento de $13/8$ a $21/8$ registra la pérdida por
$d^{-1}$. Es una contribución del despeje del coeficiente, distinta
de los errores que ya estaban incluidos en $\mathcal B$.
La prueba coincide con el núcleo numérico de A, Proposición~10.2,
pp.~48--49, y de B, Lema~6.1, pp.~65--66.

\subsection{De la selección de caracteres a una desigualdad escalar}

La Nota~\ref{sec:caracteres} retiene $4o$ caracteres: los $o$
que forman \eqref{eq:par-hpsi} y $3o$ auxiliares. Denotemos por
$\Sigma_{\mathrm{aux}}(x)$ una cota no negativa para la suma de
los valores absolutos de sus contribuciones auxiliares. Escribamos
$\rho$ para el error de los coeficientes del selector, $P_x$ para
una cota de las sumas retenidas y $V_x$ para el error producido al
cambiar un argumento $y$ por $x$ cuando $|y-x|\le1$.
El coste de reemplazar los pesos ya se incorpora al término
$4/(dx)$ que aparece abajo. Estas cantidades son las de la
extracción de A, Proposición~4.3; no son la suma cohomológica
$V(E,h')$ de la Nota~\ref{sec:cohomologia}.

Bajo las hipótesis sobre pesos, cotas analíticas y aproximación del
selector de A, Proposición~4.3, recibimos la desigualdad
\begin{equation}
 H_\psi(x)\le
 \Sigma_{\mathrm{aux}}(x)+F_x+
 \frac{2E}{Qd}+\frac{4}{dx},
 \qquad F_x=4o(V_x+\rho P_x).
 \label{eq:par-selector}
\end{equation}
Aquí $E$ acota el exceso para cada traslación del conjunto utilizado.
Para pasar de esa cota puntual al promedio se emplean sus coeficientes
no negativos y el mismo sistema de datos. Una media tomada con otros
pesos necesitaría su propia justificación.

\begin{lema}[Separación de los dos errores de extracción]
\label{lem:par-selector}
Supongamos \eqref{eq:par-selector}, $E\le x^{C\varepsilon}\mathcal B$ y
\begin{equation}
 \Sigma_{\mathrm{aux}}(x)\le4o\,x^{1/2+\varepsilon},
 \qquad F_x\le4x^{1/2+\varepsilon}.
 \label{eq:par-auxiliares}
\end{equation}
Definamos
\begin{equation}
 Z=H_\psi(x)-4o\,x^{1/2+\varepsilon}.
 \label{eq:par-z}
\end{equation}
Entonces
\begin{equation}
 \frac{Qd}{2}\left(Z-4x^{1/2+\varepsilon}\right)
 \le x^{C\varepsilon}\mathcal B+\frac{2Q}{x}.
 \label{eq:par-escalar}
\end{equation}
\end{lema}

\begin{proof}
Sustituir \eqref{eq:par-auxiliares} en
\eqref{eq:par-selector} y restar los dos términos explícitos da
\[
 Z-4x^{1/2+\varepsilon}
 \le \frac{2E}{Qd}+\frac{4}{dx}.
\]
Multiplicamos por $Qd/2>0$ y utilizamos la cota de $E$.
\end{proof}

La desigualdad \eqref{eq:par-escalar} es la hipótesis de B,
Lema~6.1, con nuestra notación para $H_\psi$. Los dos desplazamientos
son distintos: $4o\,x^{1/2+\varepsilon}$ corresponde a los caracteres
auxiliares, mientras que $4x^{1/2+\varepsilon}$ controla los demás
errores de extracción. El coeficiente $4o$ de la cota auxiliar
incorpora holgura; el número de caracteres auxiliares sigue siendo
$3o$.

El lema explica la deducción entre desigualdades. Las cotas
\eqref{eq:par-auxiliares} requieren las estimaciones de los primos
auxiliares y del selector. Una cota cruda del tipo
$P_x=O_\psi(x\log x)$ solo resulta útil en \eqref{eq:par-selector}
cuando se combina con la pequeñez de $\rho$. Por ejemplo,
$\rho=O_\psi(x^{-1+\varepsilon})$ da
$\rho P_x=O_\psi(x^\varepsilon\log x)$; para $\varepsilon$ fijo,
este error es menor que $x^{1/2+\varepsilon}$ después de un umbral.
Los factores implícitos deben permanecer fijos respecto de $x$.

\subsection{El exponente final y la absorción de constantes}

\begin{proposicion}[Lema numérico de cierre]
\label{prop:par-cierre}
Sea $Z\in\R$ que satisface \eqref{eq:par-escalar}.
Bajo \eqref{eq:par-rangos} se tiene
\begin{equation}
 Z\le26x^{1/2+(C+21/8)\varepsilon}.
 \label{eq:par-veintiseis}
\end{equation}
Si $26\le x^{3\varepsilon/8}$, entonces
$Z\le x^{1/2+(C+3)\varepsilon}$.
Si además $Z$ está dado por \eqref{eq:par-z} y
$1+4o\le x^\varepsilon$, entonces
\begin{equation}
 H_\psi(x)\le x^{1/2+(C+4)\varepsilon}.
 \label{eq:par-cierre}
\end{equation}
\end{proposicion}

\begin{proof}
Dividimos \eqref{eq:par-escalar} por $Q/2$ y aplicamos
\eqref{eq:par-diez}:
\[
 d\left(Z-4x^{1/2+\varepsilon}\right)
 \le20x^{1/2+(C+13/8)\varepsilon}+\frac4x
 \le22x^{1/2+(C+13/8)\varepsilon}.
\]
La última desigualdad usa $x\ge3$ y
$x^{1/2+(C+13/8)\varepsilon}\ge1$.
Como $d^{-1}\le x^\varepsilon$, resulta
\[
 Z\le4x^{1/2+\varepsilon}
       +22x^{1/2+(C+21/8)\varepsilon}
 \le26x^{1/2+(C+21/8)\varepsilon}.
\]
La comparación de exponentes usa $C\ge0$. No se ha supuesto que
$Z$ o $H_\psi$ sean no negativos.

La identidad $21/8+3/8=3$ muestra que
$26\le x^{3\varepsilon/8}$ absorbe el factor $26$.
Por \eqref{eq:par-z},
\[
 H_\psi(x)\le x^{1/2+(C+3)\varepsilon}
                  +4o\,x^{1/2+\varepsilon}
 \le (1+4o)x^{1/2+(C+3)\varepsilon}.
\]
Usar $1+4o\le x^\varepsilon$ prueba \eqref{eq:par-cierre}.
\end{proof}

Un umbral explícito suficiente para las dos absorciones es
\begin{equation}
 X_{\mathrm{abs}}(\psi,\varepsilon)=
 \max\left\{3,\ 26^{8/(3\varepsilon)},\
                      (1+4o)^{1/\varepsilon}\right\}.
 \label{eq:par-umbral}
\end{equation}
Debe combinarse con los umbrales de las estimaciones recibidas.
Su tamaño puede ser enorme cuando $\varepsilon$ es pequeño.
La conclusión de existencia de ventanas arbitrariamente lejanas
permite imponerlo sin exigir un valor práctico de $Y$.

A especializa los exponentes de divisibilidad y pérdida a
$e_D=2$ y $e_L=0$, y utiliza $C=e_D+e_L+3=5$.
En esa especialización \eqref{eq:par-cierre} toma la forma
$H_\psi(x)\le x^{1/2+9\varepsilon}$.
Conservar $C$ en la prueba permite ver qué uniformidad necesita
la conclusión; el valor $5$ requiere la especialización concreta
de las entradas de la fuente.

\subsection{Realización uniforme y orden de las elecciones}

La siguiente hipótesis expresa el contenido que B reúne en su
Proposición~6.3. Fijado $\psi$, existen $C\ge0$ y
$\varepsilon_0>0$ tales que, para cada
$0<\varepsilon<\varepsilon_0$ y cada $A\ge3$, hay $Y\ge A$
con esta propiedad: para todo $x\in[Y,Y^{3/2}]$ se pueden elegir
datos que satisfacen \eqref{eq:par-rangos} y
\eqref{eq:par-escalar}, con $Z$ dado por \eqref{eq:par-z}.
Llamaremos a esta afirmación $\mathsf U_\psi$.

Su estructura es
\begin{equation}
 \exists C,\varepsilon_0\
 \forall\varepsilon\in(0,\varepsilon_0)\
 \forall A\ge3\
 \exists Y\ge A\
 \forall x\in[Y,Y^{3/2}]\
 \exists\text{ datos compatibles}.
 \label{eq:par-cuantificadores}
\end{equation}
Los valores de $C$ y $\varepsilon_0$ pueden depender de $\psi$.
El primero no depende de $\varepsilon$, de $Y$ ni de $x$.
Los últimos datos pueden depender de todos los parámetros que los
preceden. La existencia de números que cumplen los rangos de
\eqref{eq:par-rangos} es solo una parte de esta hipótesis:
también debe existir la construcción que justifica
\eqref{eq:par-escalar}.

\begin{proposicion}[De la realización uniforme a las ventanas]
\label{prop:par-ventanas}
La hipótesis $\mathsf U_\psi$ implica \eqref{eq:par-meta}.
\end{proposicion}

\begin{proof}
Fijemos $\eta>0$ y $A\ge3$. Elegimos primero los valores de
$C$ y $\varepsilon_0$ suministrados por $\mathsf U_\psi$,
y después un número
\[
 0<\varepsilon<
       \min\left\{\varepsilon_0,\frac{\eta}{C+4}\right\}.
\]
Aplicamos $\mathsf U_\psi$ con el umbral
$A_*=\max\{A,X_{\mathrm{abs}}(\psi,\varepsilon)\}$.
Obtenemos un solo $Y\ge A_*$ para el cual los datos existen en
cada punto $x\in[Y,Y^{3/2}]$. La Proposición~\ref{prop:par-cierre}
se aplica en todos esos puntos con el mismo $C$ y el mismo
$\varepsilon$, y da
\[
 H_\psi(x)\le x^{1/2+(C+4)\varepsilon}
              \le x^{1/2+\eta}.
\]
El último paso usa $x\ge1$.
\end{proof}

La prueba muestra la función de cada cuantificador. Si se obtuviera
un $C(\varepsilon)$ arbitrario, no se seguiría que
$(C(\varepsilon)+4)\varepsilon$ pudiera hacerse pequeño.
Si el umbral dependiera libremente de $x$, tampoco se obtendría una
cota en el intervalo elegido. A, §11, conserva una familia de datos,
sus constantes y un umbral común antes de escoger la ventana final.

\subsection{Una misma pareja de primos en la ventana}

La Proposición~\ref{prop:cohom-ventana-aritmetica} proporciona,
bajo sus entradas, una ventana cohomológica. Para combinarla con
las estimaciones de los caracteres auxiliares se necesita la misma
pareja de primos. A, §11, pp.~49--50, organiza las dos entradas
con el siguiente orden. La fuente C desarrolla el mecanismo de las
dos cotas de torre que intervienen en la entrada cohomológica;
véase el complemento~\ref{subsec:cohom-fuente-c}. Su estudio conserva
la necesidad de una cota uniforme de generadores y de la estimación
en la capa fija (C, §7, pp.~11--12) \cite{fuenteC}.
El ensamblaje de esta nota exige además la selección analítica
de la misma pareja.
La Nota~\ref{sec:auxiliares} reconstruye esa selección bajo las
entradas analíticas que identifica.

Fijados $\psi$ y $0<\varepsilon<1/52$, la selección analítica
afirma que, para cada $J\ge1$ fijado, todo $T$ suficientemente grande
admite una pareja $q_{\mathrm{aux}},d_{\mathrm{aux}}$ que satisface
las condiciones analíticas en $[T,T^J]$. La pareja seleccionada
cumple además las congruencias y tamaños exigidos por la entrada
cohomológica. Estas son propiedades del teorema de selección que
recibimos; no se deducen de la mera existencia de dos primos.

Por su parte, fijemos un selector diádico $h:[1,\infty)\to\Z_{\ge0}$
tal que
\begin{equation}
 y^{2+8\varepsilon}\le2^{h(y)}
                  \le2y^{2+8\varepsilon}.
 \label{eq:par-selector-diadico}
\end{equation}
Se puede tomar $h(y)=\lceil(2+8\varepsilon)\log_2 y\rceil$.
La entrada cohomológica elige $J\ge1$ y $X_0$ antes de la pareja:
para todo $X\ge X_0$ y toda pareja admisible a escala $X^{1/J}$,
existe $W\ge X^{1/J}$ con $W^2\le X$ tal que
\begin{equation}
 \mathcal V\bigl(E_{h(y),q_{\mathrm{aux}},d_{\mathrm{aux}}},
                         \lfloor y^{401}\rfloor\bigr)
       \le\exp(y^\varepsilon/2)
 \qquad(W\le y\le W^2).
 \label{eq:par-ventana-cohom}
\end{equation}
La estimación vale para todo cuerpo auxiliar que realiza la
descripción de cuerpo fijo requerida por esa entrada. Usamos
$\mathcal V$ con la finitud y sumabilidad ya separadas en la
Nota~\ref{sec:cohomologia}. Su monotonía en el corte fue demostrada
en el Lema~\ref{lem:cohom-presupuesto}.

El orden permite elegir primero $J$, después aplicar con ese $J$
la selección analítica y finalmente obtener $W$ para la pareja
seleccionada. Si el teorema cohomológico solo permitiera escoger
otra pareja, este argumento no conectaría las dos estimaciones.

\begin{proposicion}[Transporte a una ventana común]
\label{prop:par-ventana-comun}
Recibidas las dos entradas anteriores, sea $c\ge1$ fijo. Existen
$J\ge1$ y valores arbitrariamente grandes de $Y$ para los cuales
una misma pareja $q_{\mathrm{aux}},d_{\mathrm{aux}}$ y una misma
escala $T$ sirven para todos los $x\in[Y,Y^{3/2}]$.
En cada punto puede elegirse un entero $e\ge0$ tal que
\begin{equation}
 \begin{gathered}
 x^{1/J}\le T\le x,\qquad
 x^{2+8\varepsilon}\le c\,2^e\le2x^{2+8\varepsilon},\\
 \mathcal V\bigl(E_{e,q_{\mathrm{aux}},d_{\mathrm{aux}}},
                            \lfloor x^{400}\rfloor\bigr)
      \le\exp(x^\varepsilon/2).
 \end{gathered}
 \label{eq:par-ventana-comun}
\end{equation}
Las condiciones analíticas de la pareja permanecen válidas en $[T,x]$.
\end{proposicion}

\begin{proof}
Tomamos $J,X_0$ de la entrada cohomológica y el umbral $T_0$
de la selección analítica para ese mismo $J$. Pongamos
\[
 \lambda=c^{1/(2+8\varepsilon)}\ge1.
\]
El símbolo $\lambda$ evita confundir este factor real con el
carácter cuadrático $\kappa$ de \eqref{eq:par-periodo}.
Dado un umbral deseado $A$, elegimos
\[
 M\ge\max\{A,T_0,1,\lambda^3,\lambda^{400}\},
 \quad X\ge\max\{X_0,M^J\},\quad T=X^{1/J}.
\]
La selección analítica proporciona una pareja en $[T,X]$.
La entrada cohomológica, aplicada a esa pareja, da $W\ge T$
con $W^2\le X$ y \eqref{eq:par-ventana-cohom}.
Definimos $Y=\lambda W$.

Sea $x\in[Y,Y^{3/2}]$ y pongamos $y=x/\lambda$.
Como $W\ge\lambda^3$, se tiene
$(\lambda W)^{3/2}\le W^2$. Por tanto,
\begin{equation}
 T\le W\le y\le x\le W^2\le X.
 \label{eq:par-inclusiones}
\end{equation}
Las condiciones analíticas se restringen de $[T,X]$ a $[T,x]$.
Además, $x\le X$ implica $x^{1/J}\le X^{1/J}=T$.

Tomamos $e=h(y)$. La igualdad
$\lambda^{2+8\varepsilon}=c$ transforma
\eqref{eq:par-selector-diadico} en el rango de $c\,2^e$
de \eqref{eq:par-ventana-comun}. Para el corte de torcidos,
$y\ge W\ge\lambda^{400}$ da
\begin{equation}
 x^{400}=\lambda^{400}y^{400}\le y^{401},
 \qquad
 \lfloor x^{400}\rfloor\le\lfloor y^{401}\rfloor.
 \label{eq:par-corte}
\end{equation}
La monotonía de $\mathcal V$ permite disminuir el corte de
\eqref{eq:par-ventana-cohom}. Puesto que $y\le x$,
también se tiene $\exp(y^\varepsilon/2)\le\exp(x^\varepsilon/2)$.
Esto prueba la última cota de \eqref{eq:par-ventana-comun}.
La pareja y $T$ se eligieron antes de $x$, y
$Y\ge W\ge M\ge A$.
\end{proof}

El exponente $401$ reserva una potencia para absorber el factor fijo
de reescala; la construcción final utiliza el corte $400$.
La potencia $3/2$ de la ventana deja espacio dentro de $[W,W^2]$.
La meta cohomológica sigue siendo
$\exp(x^\varepsilon/2)=\sqrt{\exp(x^\varepsilon)}$.
La expresión $\exp(x^{\varepsilon/2})$ impondría una cota diferente.

\subsection{Un máximo finito de umbrales para una sola familia}

Una pareja común resuelve parte de la compatibilidad. Las aplicaciones
integrales, los marcos, los rangos y las métricas deben pertenecer
también a una misma familia. El principio de cierre que A usa en
la Proposición~11.2 tiene una prueba lógica breve, pero sus hipótesis
conservan toda la carga de las construcciones.

\begin{lema}[Reunión de condiciones eventuales]
\label{lem:par-umbral-comun}
Fijado $\varepsilon$, sea $\operatorname{Adm}_\varepsilon(K,\mathcal F)$,
con $K\in\R_{\ge0}$, una condición de calidad para familias,
monótona en el siguiente
sentido:
\[
 K'\ge K,\quad
 \operatorname{Adm}_\varepsilon(K',\mathcal F)
 \quad\Longrightarrow\quad
 \operatorname{Adm}_\varepsilon(K,\mathcal F).
\]
Supongamos que para cada $K$ existe una familia
$\mathcal F\in\mathcal C'$ que satisface esa condición.
Sean $P_1,\ldots,P_s$, con $s\ge1$, un número finito de propiedades. Si para cada
$i$ hay $K_i$ tal que toda familia de $\mathcal C'$ con calidad
$K_i$ tiene $P_i$, existe una sola familia de $\mathcal C'$
que satisface todas las propiedades.
\end{lema}

\begin{proof}
Tomamos $K_*=\max\{K_1,\ldots,K_s\}$ y elegimos una familia
$\mathcal F\in\mathcal C'$ con
$\operatorname{Adm}_\varepsilon(K_*,\mathcal F)$.
La monotonía da
$\operatorname{Adm}_\varepsilon(K_i,\mathcal F)$ para cada $i$,
de donde se sigue $P_i(\mathcal F)$.
\end{proof}

Si cada propiedad exige además
$0<\varepsilon<\varepsilon_i$, se fija primero
$\varepsilon<\min_i\varepsilon_i$ y después se obtienen los
$K_i(\varepsilon)$. Si el umbral es una función de otros parámetros
fijos, el máximo se toma punto a punto. Se necesita que la existencia
de familias esté disponible para ese umbral combinado.

La afirmación $\forall i\,\exists\mathcal F_i\,P_i(\mathcal F_i)$
no contiene este principio: una propiedad podría requerir un índice
par y otra uno impar. Las construcciones de A suministran propiedades
de \emph{todas} las familias suficientemente admisibles de una clase
común, y el teorema de existencia se aplica después de reunir
sus umbrales.

\subsection{Las entradas que conserva el ensamblaje geométrico}

En A, Proposición~11.2, la clase de existencia
$\mathcal C'$ está incluida en una clase geométrica $\mathcal C$
y en una clase aritmética $\mathcal E$. La inclusión
\[
 \mathcal C'\subseteq\mathcal C\cap\mathcal E
\]
permite usar ambas colecciones de propiedades para la misma familia.
La existencia debe aportar también los denominadores suaves,
los retorcimientos compatibles, una constante $C_H\ge0$ y el
testigo de evaluación \texttt{ChosenE}, junto con los controles de
tamaño \texttt{TuSize} y \texttt{DGamSize}. Estas son entradas
adicionales de la construcción, no consecuencias de la notación
de clases.

Los nueve grupos de propiedades eventuales de la fuente tienen
los siguientes cometidos. Los nombres permiten localizar las
interfaces; la tabla no afirma una revisión completa de sus pruebas.

\begin{center}
\small
\begin{tabularx}{\textwidth}{@{}p{0.29\textwidth}Y@{}}
\toprule
\textbf{Propiedad de A} & \textbf{Entrada que debe proporcionar}\\
\midrule
\texttt{LinkN} &
Control del defecto integral mediante longitudes cohomológicas
y errores de tamaño.\\
\texttt{LinkEta} &
Condición de cubrimiento local compatible con los torcidos elegidos.\\
\texttt{LinkIRR} &
Cota cuantitativa de la proporción de filas eliminadas y de la
fracción de rango retenida.\\
\texttt{SmallLevelRankEq} &
Identificación del rango después del paso a nivel pequeño.\\
\texttt{LinkDIVt} &
Mapas integrales de evaluación y comparaciones de altura, periodo,
nivel y pendiente.\\
\texttt{NuFullF}, \texttt{MuRank} &
Plenitud de los retículos y positividad de los rangos necesarios.\\
\texttt{LinkG} &
Ganancia en lugares finitos, costes de índices y comparación de
órbitas de Frobenius servidas.\\
\texttt{LinkS} &
Estimaciones espectrales y de marco aproximado.\\
\texttt{LinkGLOBAL} &
Comparación global de pendientes que reúne esas entradas.\\
\bottomrule
\end{tabularx}
\end{center}

El rango positivo permite dividir por un rango. La cota cuantitativa
de \texttt{LinkIRR} controla, además, cuánto puede costar esa división.
Son hipótesis diferentes; la positividad no sustituye a una
proporción retenida cercana a uno.

La clase aritmética $\mathcal E$ conserva dos propiedades: la
estimación cohomológica eventual para cada dato pertinente y la
existencia de intervalos completos $[Y,Y^{3/2}]$ arbitrariamente
lejanos dentro del conjunto de escalas admisibles. La reunión finita
se hace dentro de esa clase, de modo que la familia elegida mantiene
las ventanas.

Se distinguen dos umbrales. El umbral de calidad $K_*$ se fija antes
de invocar la existencia de $\mathcal F_\varepsilon$.
Después se fijan las constantes de comparación de esa familia y
un umbral de escala $x_0$. Solo entonces se elige
$[Y,Y^{3/2}]$ con $Y\ge x_0$ dentro de sus escalas admisibles.
La sucesión de dependencias puede escribirse como
\begin{equation}
 \begin{gathered}
 \psi,\ (e_D,e_L)
 \ ;\quad \varepsilon
 \ ;\quad K_*
 \ ;\quad \mathcal F_\varepsilon;\\
 \text{constantes y }x_0
 \ ;\quad Y
 \ ;\quad x\text{ y sus datos}.
 \end{gathered}
 \label{eq:par-orden-familia}
\end{equation}
Aquí $C=e_D+e_L+3$ está fijado antes de $\varepsilon$.
Los prefactores y $x_0$ pueden depender de la familia y de
$\varepsilon$; su uniformidad respecto de $x$ permite escoger
la ventana después.

\subsection{El enlace con el criterio unilateral}

La Proposición~\ref{prop:par-ventanas} recibe
$\mathsf U_\psi$ como entrada. Las
Proposiciones~\ref{prop:par-ventana-comun} y el
Lema~\ref{lem:par-umbral-comun} reconstruyen dos operaciones
de su realización: conservar una pareja durante una ventana y
reunir propiedades eventuales en una familia. Las estimaciones
especializadas que alimentan esas operaciones siguen siendo
necesarias.

Si se dispone de esas entradas para cada carácter primitivo, el
paso analítico final puede formularse sin añadir una cota entre
ventanas. Supongamos que $\rho$ es un cero de $L(s,\psi)$ con
$\beta=\Ree\rho>1/2$. Elegimos
$0<\eta<\min\{\beta-1/2,1/2\}$ y ponemos
$\theta=1/2+\eta\in(1/2,1)$. La ventana
\eqref{eq:par-meta} se convierte, al escribir $x=e^u$ y
$U=\log Y$, en
\begin{equation}
 H_\psi(e^u)\le e^{\theta u}
 \qquad(U\le u\le3U/2).
 \label{eq:par-log-ventana}
\end{equation}
Es la ventana de la Nota~\ref{sec:ventanas} con $\delta=1/2$.
El mayorante $R_U(u)=e^{\theta u}$ satisface
\[
 \int_U^{3U/2}R_U(u)e^{-\theta u}\dd u
       =\frac U2.
\]
Para todo $\eta'>0$ fijo, $U/2\le e^{\eta'U}$ cuando $U$
es suficientemente grande. Las ventanas arbitrariamente lejanas
permiten imponer ese umbral. El criterio unilateral excluye
entonces el polo asociado a $\rho$, pues $\beta>\theta$ y
el residuo no se cancela entre los coeficientes positivos de
\eqref{eq:par-hpsi}. La ecuación funcional proporciona la reducción
del otro lado de la franja ya expuesta en la introducción.

Esta es una deducción condicionada a la realización uniforme.
La cota obtenida concierne a $H_\psi$, no a $|H_\psi|$, y se
afirma en las ventanas seleccionadas. Las Notas~\ref{sec:sumas}
y~\ref{sec:ventanas} explican por qué esa información basta
para el argumento analítico.

\subsection{Concordancia localizada con Lean y alcance}

El cotejo de este bloque utiliza la revisión fijada
\texttt{7cbf7a0} del espejo indicado en la referencia~\cite{lean}.
Comprende 28 archivos: 26 leídos completos y dos de manera selectiva,
con 4497 líneas leídas. Se recuperó además un archivo para localizar
el recorrido, excluido de las conclusiones. El registro conserva
las rutas, los intervalos y los hashes de los archivos consultados.
Los siguientes enlaces apuntan a esa revisión inmutable.

\paragraph{El presupuesto del resto.}
\href{https://github.com/CocoSgt/riemann-hypothesis-lean/blob/7cbf7a0c32fca1badf05931879a37c3f8a2b03b4/src/Formalism/SizeWindowBrackets/ErrorFormsWeakWindowsFalse.lean}{\texttt{Error\allowbreak{}Forms\allowbreak{}Weak\allowbreak{}Windows\allowbreak{}False}}
define \texttt{ErrWindows} con una condición de traza que conviene
mantener explícita:
\begin{equation}
 (\bar\eta+4\sqrt{\bar e})x^{c_{\mathrm{HH}}}b
       \le Kx^\varepsilon(n_Ta).
 \label{eq:par-traza-lean}
\end{equation}
Aquí $\bar\eta,\bar e$ miden errores del marco; $n_Ta$ es la
escala de referencia de ese presupuesto, y $c_{\mathrm{HH}}$
controla el crecimiento del factor geométrico. Estos símbolos
pertenecen a esa interfaz. La condición \eqref{eq:par-traza-lean}
controla el producto completo: las cotas
$\bar\eta\le1$ y $\bar e\le1$ dejarían sin controlar la
amplificación por $x^{c_{\mathrm{HH}}}b$.

En \href{https://github.com/CocoSgt/riemann-hypothesis-lean/blob/7cbf7a0c32fca1badf05931879a37c3f8a2b03b4/src/Formalism/SizeWindowBrackets/ChainSlopeWindows/SubThreeMulLogTwoLe.lean}{\texttt{SubThreeMulLogTwoLe}},
\texttt{error\_sizes} prueba, bajo \texttt{ErrWindows}
y la condición de escala \texttt{LargeX},
\[
 \mathrm{errorTermWith}\le
 x^{(e_{\mathrm{Div}}+e_{\mathrm{Loss}}+3)\varepsilon}
 \left[n_Ta\left(1+\frac{m^2}{k}\right)+m(B_0+N_c)\right].
\]
La constante que se absorbe está fijada antes de aumentar $x$:
\[
 K_{\mathrm{big}}=200(1+K)^2(1+c_{\mathrm{HH}})(1+e_K),
 \qquad K_{\mathrm{big}}\le x^\varepsilon.
\]
El mismo archivo prueba \texttt{error\_sizes\_weak\_false},
un contraejemplo numérico a la interfaz que sustituye
\eqref{eq:par-traza-lean} por las dos cotas débiles.
Su alcance es esa sustitución de hipótesis; no constituye un
contraejemplo a la construcción aritmética ni a RH.
\href{https://github.com/CocoSgt/riemann-hypothesis-lean/blob/7cbf7a0c32fca1badf05931879a37c3f8a2b03b4/src/Development/AutomorphicArithmetic/DedekindLFunctions/LandauTheorems/ChTop.lean}{\texttt{ChTop}} transporta el presupuesto al exceso
\texttt{gapA} de cada traslación, conservando
\texttt{IdentU}, \texttt{SizeData}$_2$, \texttt{BNBoundsU},
\texttt{ErrWindows} y \texttt{LargeX}.

\paragraph{El umbral y la combinación retenida.}
En \href{https://github.com/CocoSgt/riemann-hypothesis-lean/blob/7cbf7a0c32fca1badf05931879a37c3f8a2b03b4/src/Development/AutomorphicArithmetic/DedekindLFunctions/LandauTheorems/ChCore.lean}{\texttt{ChCore}},
\texttt{ch\_core} fija $C,\varepsilon,C_E,K_o,c,M$ antes
de elegir $x_1$. El grupo, el host, los conjuntos, los pesos y
los demás datos aparecen después de ese umbral. Bajo las
condiciones de tamaños, escalas, caracteres y exceso, obtiene
\[
 G_{m,A}(x)\le x^{1/2+(C+3)\varepsilon}+\Sigma_A(x).
\]
\href{https://github.com/CocoSgt/riemann-hypothesis-lean/blob/7cbf7a0c32fca1badf05931879a37c3f8a2b03b4/src/Development/AutomorphicArithmetic/PrimitiveCosets/SigALe.lean}{\texttt{SigALe}} recibe la cota analítica
\texttt{EAtH.bdG} y demuestra
$\Sigma_A(x)\le4o\,x^{1/2+\varepsilon}$.
\href{https://github.com/CocoSgt/riemann-hypothesis-lean/blob/7cbf7a0c32fca1badf05931879a37c3f8a2b03b4/src/Development/AutomorphicArithmetic/PrimitiveCosets/CoreGm.lean}{\texttt{CoreGm}} usa $C=e_D+e_L+3$
y $1+4o\le x^\varepsilon$ para concluir
\begin{equation}
 \Lambda.\mathrm{Gm}(x)
       \le x^{1/2+(e_D+e_L+7)\varepsilon}.
 \label{eq:par-exponente-lean}
\end{equation}
Con las identificaciones del nivel cíclico
$\Lambda=\mathrm{Level.cyc}(\psi)$,
$\Lambda.\mathrm{Gm}$ corresponde a $H_\psi$; no designa
la suma individual $G_\chi$. El exponente coincide con
\eqref{eq:par-cierre}, pues $C+4=e_D+e_L+7$.
Otras interfaces refinadas del recorrido solo-zeta muestran
$1/2+6\varepsilon$ al especializar $e_D=2,e_L=0$;
ese exponente pertenece a otra formulación y no reemplaza
\eqref{eq:par-exponente-lean}.

\paragraph{Existencia y estimación sobre una misma familia.}
\href{https://github.com/CocoSgt/riemann-hypothesis-lean/blob/7cbf7a0c32fca1badf05931879a37c3f8a2b03b4/src/Formalism/SharpRestriction/RamificationFrobenius/TotientGaloisClasses/TopWin.lean}{\texttt{TopWin}} distingue
\texttt{ExWin}, de la forma $\forall K\,\exists\mathcal F$,
y \texttt{ChainWin}, de la forma
$\exists K\,\forall\mathcal F$ compatible.
El teorema \texttt{tiedWin\_of\_ex} elige primero el
$K$ que proporciona la estimación y después la familia que
lo realiza. Los exponentes $e_D,e_L$ se fijan antes de
$\varepsilon$. \href{https://github.com/CocoSgt/riemann-hypothesis-lean/blob/7cbf7a0c32fca1badf05931879a37c3f8a2b03b4/src/Development/AutomorphicArithmetic/CharacterBoxes/DesWin.lean}{\texttt{DesWin}}
consume \texttt{DesignExists}, nueve hipótesis
\texttt{Closed}, la entrada cohomológica \texttt{HypV}
y las compatibilidades entre clases.
Estas son pruebas de ensamblaje bajo entradas expresas.

\paragraph{La selección analítica de los auxiliares.}
\href{https://github.com/CocoSgt/riemann-hypothesis-lean/blob/7cbf7a0c32fca1badf05931879a37c3f8a2b03b4/src/Development/AutomorphicArithmetic/CharacterBoxes/PairGood.lean}{\texttt{PairGood}} define \texttt{PairGood} como
un paquete de propiedades aritméticas, cotas de torcidos en
una ventana completa y cajas libres de ceros. La proposición
\texttt{PairEx} afirma la existencia de esos pares, pero ese
archivo no contiene una prueba de su realización.
\href{https://github.com/CocoSgt/riemann-hypothesis-lean/blob/7cbf7a0c32fca1badf05931879a37c3f8a2b03b4/src/Development/AutomorphicArithmetic/CharacterBoxes/PairRows.lean}{\texttt{PairRows}} transporta el paquete a las entradas
del host. \href{https://github.com/CocoSgt/riemann-hypothesis-lean/blob/7cbf7a0c32fca1badf05931879a37c3f8a2b03b4/src/Development/AutomorphicArithmetic/CharacterBoxes/WinPairs.lean}{\texttt{WinPairs}} recibe
\texttt{PairEx} y \texttt{VWinAt} de profundidad 401;
su reescalamiento utiliza las condiciones correspondientes a
$W\ge\lambda^3$ y $W\ge\lambda^{400}$ de
la Proposición~\ref{prop:par-ventana-comun}.
La lectura confirma las inclusiones de intervalos y conserva
la selección analítica como entrada.

\paragraph{El paso analítico final.}
\href{https://github.com/CocoSgt/riemann-hypothesis-lean/blob/7cbf7a0c32fca1badf05931879a37c3f8a2b03b4/src/Development/AutomorphicArithmetic/DedekindLFunctions/LandauTheorems/WindowHypOfRpow.lean}{\texttt{WindowHypOfRpow}} convierte las cotas puntuales
en una condición integral mediante
$R(y)=|C_0|y^\theta$, con $C_0$ fijo. Demuestra la
integrabilidad y la cuenta
\[
 \int_Y^{Y^{1+\delta}}R(y)y^{-\theta-1}\dd y
       =|C_0|\delta\log Y\le Y^{\eta'}
\]
para $Y$ suficientemente grande, dado $\eta'>0$.
\href{https://github.com/CocoSgt/riemann-hypothesis-lean/blob/7cbf7a0c32fca1badf05931879a37c3f8a2b03b4/src/Development/AutomorphicArithmetic/DedekindLFunctions/LandauTheorems/WindowHypUOfWindowHyp.lean}{\texttt{WindowHypUOfWindowHyp}} realiza $y=e^u$
y conserva integrabilidad y cuantificadores.
\href{https://github.com/CocoSgt/riemann-hypothesis-lean/blob/7cbf7a0c32fca1badf05931879a37c3f8a2b03b4/src/Development/AutomorphicArithmetic/DedekindLFunctions/LandauTheorems/landauWindow.lean}{\texttt{landauWindow}} excluye los polos con
parte real mayor que $\theta$, bajo \texttt{PoleData},
$0<\theta<1$, $\delta>0$ y \texttt{WindowHypU}.
Son los pasos analíticos utilizados en el apartado anterior.

Las dos lecturas selectivas son
\texttt{GrowthBracketBoundOfSizeWindows}, líneas 100--165,
y \texttt{PChain}, líneas 1--155. Las definiciones de
proposiciones y los campos de estructuras no se cuentan como
pruebas de la existencia de sus datos. No se ejecutó ni instaló
Lean, no se auditó la clausura completa de imports o axiomas
ni se comparó todo el espejo con la distribución original.
Este cotejo localizado y las deducciones de la nota no certifican
el argumento global.

\subsection{Ejercicios con soluciones}

\begin{ejercicio}\label{ej:par-equilibrio}
Sean $A_0,B_0,x>0$. Determine el mínimo, para $m>0$, de
$A_0x/m+B_0m$. Explique qué cambia si solo se puede elegir
$m$ entero y por qué el equilibrio de este ejercicio no prueba
por sí solo la existencia de una torre theta admisible.
\end{ejercicio}

\paragraph{Solución del Ejercicio~\ref{ej:par-equilibrio}.}
Se tiene
\[
 \frac{A_0x}{m}+B_0m-2\sqrt{A_0B_0x}
 =\left(\sqrt{B_0m}-\sqrt{A_0x/m}\right)^2.
\]
El mínimo es $2\sqrt{A_0B_0x}$ y se alcanza en
$m_*=\sqrt{A_0x/B_0}$. Para $m$ entero positivo se compara
entre los enteros vecinos permitidos, usando la monotonía de
la función a cada lado de $m_*$. Si $m_*<1$, el mínimo entero
está en $m=1$; si $m_*\ge1$, basta comparar
$\lfloor m_*\rfloor$ y $\lceil m_*\rceil$.
Cuando $A_0,B_0$ son fijos y $x$ crece, el redondeo es
despreciable en relación con $m_*$.
El ejercicio optimiza una función real; la torre debe cumplir
también las condiciones de rango, integridad, pesos y compatibilidad
de la construcción.

\begin{ejercicio}\label{ej:par-modelo}
Sea $s\ge2$ entero, $x=s^2$, $0<\varepsilon<1$ y $b>0$.
Considere
\[
 n=t=v_2=x,\qquad m=s,\qquad k=s^3,\qquad a_E=2xb.
\]
Compruebe \eqref{eq:par-rangos} y calcule $\mathcal B/Q$.
Indique el alcance de este ejemplo.
\end{ejercicio}

\paragraph{Solución del Ejercicio~\ref{ej:par-modelo}.}
Las dos densidades son uno y las condiciones de $t,v_2,a_E$
se cumplen. Como $m=\sqrt{x}$, también se cumple su rango.
Se tiene $k=xm=x^{3/2}$ y
\[
 n^2x^{1/2-\varepsilon}
 =x^{5/2-\varepsilon}\ge x^{3/2}
\]
porque $\varepsilon<1$ y $x\ge1$. Sustituyendo en
\eqref{eq:par-cuatro} obtenemos
\begin{equation}
 \frac{\mathcal B}{Q}
 =2\sqrt{x}+2+2x^{1/2+\varepsilon}
                         +x^{-1/2+2\varepsilon}.
 \label{eq:par-modelo}
\end{equation}
Cada término cumple el mayorante correspondiente del
Lema~\ref{lem:par-presupuesto}. El ejemplo realiza los rangos
numéricos con enteros; no proporciona el grupo de Galois, sus
posiciones, los periodos ni la desigualdad
\eqref{eq:par-escalar}.

\begin{ejercicio}\label{ej:par-cuantificadores}
\begin{enumerate}[label=\alph*)]
\item Explique por qué permitir $C(\varepsilon)=1/\varepsilon$
en una cota $H(x)\le x^{1/2+C(\varepsilon)\varepsilon}$
no permite deducir cotas con exponente arbitrariamente próximo
a $1/2$.
\item Construya dos conjuntos disjuntos de escalas que contengan
ventanas $[Y,Y^{3/2}]$ arbitrariamente lejanas.
\end{enumerate}
\end{ejercicio}

\paragraph{Solución del Ejercicio~\ref{ej:par-cuantificadores}.}
En a), el exponente de la cota es siempre $3/2$.
Por ejemplo, $H(x)=x$ satisface esa cota para $x\ge1$ y
todo $\varepsilon>0$, pero no satisface
$H(x)\le x^{1/2+\eta}$ para $x>1$ y $0<\eta<1/2$.
Esto demuestra la necesidad de controlar la dependencia de $C$.

Para b), pongamos
\[
 \mathcal W_0=\bigcup_{j\ge1}
      [e^{4^{2j}},e^{(3/2)4^{2j}}],
 \qquad
 \mathcal W_1=\bigcup_{j\ge1}
      [e^{4^{2j+1}},e^{(3/2)4^{2j+1}}].
\]
Cada intervalo tiene la forma requerida y sus extremos inferiores
tienden a infinito. En coordenadas logarítmicas son intervalos
$[4^r,(3/2)4^r]$ con $r$ de paridad distinta. El extremo
superior de cualquiera es menor que el extremo inferior del
siguiente, pues $(3/2)4^r<4^{r+1}$.
Los dos conjuntos son disjuntos. Las ventanas comunes de la
Proposición~\ref{prop:par-ventana-comun} dependen de la compatibilidad
de las entradas; no se obtienen intersectando dos conclusiones
independientes de existencia.

\begin{ejercicio}\label{ej:par-reescala}
Sea $\lambda\ge1$, $1<\alpha<2$ y $p>q>0$.
Dé condiciones sobre $W$ que garanticen, al poner
$Y=\lambda W$ y $y=x/\lambda$,
\[
 x\in[Y,Y^\alpha]\ \Longrightarrow\
 W\le y\le x\le W^2
 \quad\hbox{y}\quad x^q\le y^p.
\]
Recupere los exponentes $3$ y $400$ usados antes.
\end{ejercicio}

\paragraph{Solución del Ejercicio~\ref{ej:par-reescala}.}
La cota $x\ge Y$ da $y\ge W$ y $\lambda\ge1$ da $y\le x$.
Para asegurar $x\le W^2$ basta
\[
 (\lambda W)^\alpha\le W^2
 \quad\Longleftrightarrow\quad
 W\ge\lambda^{\alpha/(2-\alpha)}.
\]
Como $x^q=\lambda^q y^q$, para el segundo requisito basta
$y^{p-q}\ge\lambda^q$, que se cumple si
\[
 W\ge\lambda^{q/(p-q)}.
\]
Tomar $\alpha=3/2$, $p=401$ y $q=400$ produce los umbrales
$W\ge\lambda^3$ y $W\ge\lambda^{400}$. El argumento explica
las reservas cuantitativas de la elección sin atribuirles optimalidad.
